\documentclass[12pt,a4paper]{article}

\usepackage[margin=1in]{geometry}
\usepackage[T1]{fontenc}
\usepackage[utf8]{inputenc}
\usepackage{array}
\usepackage{longtable}
\usepackage{hyperref}
\usepackage{xcolor}
\usepackage{enumitem}
\usepackage{fancyhdr}
\usepackage{lastpage}
\usepackage{titlesec}

\hypersetup{
    colorlinks=true,
    linkcolor=blue,
    urlcolor=blue,
    pdftitle={May\_04\_Deniz Software Requirements Specification},
    pdfauthor={May\_04\_Deniz Project Team}
}

\setlength{\parskip}{0.6em}
\setlength{\parindent}{0pt}
\setlist[itemize]{topsep=0.2em,itemsep=0.2em}
\setlist[enumerate]{topsep=0.2em,itemsep=0.2em}
\renewcommand{\arraystretch}{1.25}
\newcolumntype{L}[1]{>{\raggedright\arraybackslash}p{#1}}

\pagestyle{fancy}
\fancyhf{}
\lhead{May\_04\_Deniz SRS}
\rhead{Whole Project}
\cfoot{\thepage\ of \pageref{LastPage}}

\titleformat{\section}{\Large\bfseries}{\thesection}{0.7em}{}
\titleformat{\subsection}{\large\bfseries}{\thesubsection}{0.7em}{}

\begin{document}

\begin{titlepage}
    \centering
    \vspace*{2cm}
    {\Huge\bfseries Software Requirements Specification\par}
    \vspace{0.6cm}
    {\Large May\_04\_Deniz Exam Monitoring and Management Platform\par}
    \vspace{1.5cm}
    {\large Whole-Project SRS\par}
    \vfill
    \begin{tabular}{ll}
        Project baseline: & \texttt{May\_04\_Deniz/} \\
        Document date: & 2026-05-12 \\
        Document version: & 1.0 \\
        Intended use: & Course review, implementation handoff, Overleaf submission \\
    \end{tabular}
    \vfill
\end{titlepage}

\tableofcontents
\newpage

\section{Introduction}

\subsection{Purpose}

This Software Requirements Specification defines the expected behavior of the
\texttt{May\_04\_Deniz} exam monitoring and management platform. It records
functional requirements, non-functional requirements, interfaces, persistence
contracts, operational constraints, and acceptance criteria for the current
implementation.

The document is written as a requirements source for reviewers, implementers,
maintainers, and future contributors. It intentionally separates what the system
shall do from how a particular function is internally implemented.

\subsection{Scope}

The system is a local-area exam platform with one authoritative server, multiple
student clients, Tk and Qt graphical interfaces, local inter-process
communication, LAN HTTP/WebSocket communication, monitoring, incident handling,
authentication validation, submission transfer, projector display, and offline
Windows installation.

\subsection{Definitions}

\begin{longtable}{L{0.22\textwidth} L{0.68\textwidth}}
\textbf{Term} & \textbf{Meaning} \\
\hline
AD & Active Directory authentication or token generation path. \\
CATS & School authentication/preflight service used before client login. \\
Dashboard & Server-side operator GUI implemented in Tk and Qt. \\
Incident & Policy-relevant event such as a blocked process, blocked title, idle state, or evidence state. \\
LAN WebSocket & Student/server runtime protocol exposed by the server at \texttt{/ws}. \\
Loopback IPC & Local-only WebSocket IPC bound to \texttt{127.0.0.1} with token authentication. \\
Projector & Read-only public display page for aggregate exam notifications. \\
SRS & Software Requirements Specification. \\
\end{longtable}

\subsection{References}

\begin{itemize}
    \item \texttt{docs/FINAL\_REPORT\_May\_04\_Deniz/02\_srs.md}
    \item \texttt{docs/FINAL\_REPORT\_May\_04\_Deniz/05\_interface\_control\_document.md}
    \item \texttt{May\_04\_Deniz/server}
    \item \texttt{May\_04\_Deniz/client}
    \item \texttt{May\_04\_Deniz/common}
    \item \texttt{May\_04\_Deniz/tests}
\end{itemize}

\section{Overall Description}

\subsection{Product Perspective}

The platform is a local exam runtime, not a cloud service. The server is the
authority for users, policy, incidents, exam state, submissions, artifacts, and
operator commands. Clients are responsible for student-side monitoring,
material extraction, local logging, replay capture, buffering, and final
submission packaging.

The architecture separates three communication domains.

\begin{longtable}{L{0.20\textwidth} L{0.35\textwidth} L{0.35\textwidth}}
\textbf{Domain} & \textbf{Purpose} & \textbf{Boundary} \\
\hline
LAN HTTP & Login, config, exam files, submissions, artifacts, projector page. & Protected routes validate session UUID; projector is public-safe. \\
LAN WebSocket & Live runtime events between server and clients. & Session UUID plus optional secured event envelope. \\
Loopback IPC & Same-machine GUI/manager communication. & Bound to loopback and token-protected, with stdio fallback. \\
\end{longtable}

\subsection{User Classes}

\begin{longtable}{L{0.22\textwidth} L{0.32\textwidth} L{0.36\textwidth}}
\textbf{User Class} & \textbf{Description} & \textbf{Primary Needs} \\
\hline
Operator & Exam supervisor using CLI or GUI. & Start and finish exams, review incidents, control users, update policy. \\
Student & Person taking the exam through the client runtime. & Authenticate, receive materials, view timer, submit work. \\
Administrator & Person preparing machines and configuration. & Configure users, auth, installer, dependencies, and network. \\
Maintainer & Developer or support engineer. & Test, debug, extend, and rebuild the system. \\
Audience & People viewing the projector page. & See large public-safe exam notifications. \\
\end{longtable}

\subsection{Operating Environment}

The primary target environment is Windows 11 x64. The software uses Python,
\texttt{aiohttp}, \texttt{psutil}, \texttt{cryptography}, \texttt{requests},
\texttt{beautifulsoup4}, and optionally \texttt{PySide6}. FFmpeg is used for
replay recording and media artifact generation.

\subsection{Constraints}

\begin{itemize}
    \item The server remains the source of truth for exam state.
    \item LAN WebSocket event schemas remain compatible between client and server.
    \item Loopback IPC is not the same protocol as LAN student/server communication.
    \item Projector output is read-only and public-safe.
    \item Temporary authentication disable is bounded and admin-managed.
    \item Manual terminal workflows remain supported through stdio fallback.
    \item The offline installer avoids global machine \texttt{site-packages}.
\end{itemize}

\section{Functional Requirements}

\subsection{Server Runtime}

\begin{longtable}{L{0.15\textwidth} L{0.47\textwidth} L{0.28\textwidth}}
\textbf{ID} & \textbf{Requirement} & \textbf{Acceptance Criteria} \\
\hline
FR-SRV-001 & The server shall start with configurable server id, host, port, broadcast interval, discovery interval, exam duration, upload limits, UI backend, IPC transport, GUI flag, reset flag, and auth secret. & Invalid startup values are rejected before runtime. \\
FR-SRV-002 & The server shall reject duplicate active servers using the same server id. & Duplicate discovery exits before binding. \\
FR-SRV-003 & The server shall expose HTTP routes for health, auth status, login, config, files, submission, artifact upload, projector, SSE, and WebSocket. & Routes are registered by the server app factory. \\
FR-SRV-004 & The server shall persist users, incidents, policy, process blacklist, process definitions, and incident rules. & Missing state files are initialized or normalized. \\
FR-SRV-005 & The server shall maintain authoritative per-user session state. & Dashboard and WebSocket state derive from the same model. \\
FR-SRV-006 & The server shall broadcast timer and session updates. & Connected clients receive periodic state. \\
FR-SRV-007 & The server shall execute admin commands from CLI and dashboard dispatch. & Exam, user, policy, auth, settings, and process commands are supported. \\
FR-SRV-008 & Global finish shall skip banned users. & Banned users remain banned and skipped count is reported. \\
FR-SRV-009 & WebSocket close reasons shall fit close-frame byte limits. & Long UTF-8 reasons are safely trimmed. \\
\end{longtable}

\subsection{Client Runtime}

\begin{longtable}{L{0.15\textwidth} L{0.47\textwidth} L{0.28\textwidth}}
\textbf{ID} & \textbf{Requirement} & \textbf{Acceptance Criteria} \\
\hline
FR-CLI-001 & The client shall start with configurable login, password, server id, direct host, port, discovery timeout, reconnect delay, recorder flag, login-check mode, UI backend, IPC transport, AD domain, and auth secret. & Invalid arguments fail before connection. \\
FR-CLI-002 & The client shall resolve its server by direct host/port or discovery. & Direct host bypasses discovery; discovery honors server id and timeout. \\
FR-CLI-003 & The client shall check health and log in before opening WebSocket. & Login returns a session UUID. \\
FR-CLI-004 & The client shall fetch exam configuration and materials after login. & Requests use the assigned session UUID. \\
FR-CLI-005 & The client shall safely extract exam materials to the student Desktop Exam folder. & Unsafe ZIP paths are rejected and unmarked user folders are preserved. \\
FR-CLI-006 & The client shall synchronize timer/submission UI with server state. & UI receives start, pause, resume, remaining time, finish, upload, and folder state. \\
FR-CLI-007 & The client shall keep local logging alive during transient WebSocket disconnects. & Local monitors and replay do not stop solely because the socket disconnects. \\
FR-CLI-008 & The client shall reconnect automatically until final submission, intentional shutdown, or process exit. & Main loop retries after configured delay. \\
FR-CLI-009 & The client shall support launcher check-login mode. & Check-login validates and exits without full runtime. \\
\end{longtable}

\subsection{Monitoring, Rules, and Incidents}

\begin{longtable}{L{0.15\textwidth} L{0.47\textwidth} L{0.28\textwidth}}
\textbf{ID} & \textbf{Requirement} & \textbf{Acceptance Criteria} \\
\hline
FR-MON-001 & The client shall collect process snapshots and detect process policy violations. & Process logs are written and matches open incidents. \\
FR-MON-002 & Executable matching shall support exact, contains, and wildcard-style entries where definitions allow it. & One entry can match related app executable variants. \\
FR-MON-003 & Focused-window titles shall be normalized before matching. & Unicode control and invisible characters are sanitized. \\
FR-MON-004 & Focused-window matching shall support legacy lists and incident rules. & Existing policy files remain valid. \\
FR-MON-005 & Saved titlebar rules shall prefer reusable contains patterns. & A pattern such as \texttt{whatsapp} matches browser title variants. \\
FR-MON-006 & Idle state shall be monitored when enabled. & Warning and critical thresholds are configurable. \\
FR-INC-001 & Incident rules shall normalize statuses unknown, whitelist, warning, and blacklist. & Common rule logic normalizes rule data. \\
FR-INC-002 & Whitelist rules shall suppress candidates before warning or blacklist rules. & Approved titles do not open incidents. \\
FR-INC-003 & Warning and blacklist rules shall attach matched rule and action metadata. & Incident payloads include rule id and configured actions. \\
FR-INC-004 & Incident rules shall persist in \texttt{data/server/incident\_rules.json}. & Load, save, import, export, and version are supported. \\
FR-INC-005 & Dashboard shall expose an Incident Rules tab. & Tk and Qt dashboards show rule management. \\
FR-INC-006 & Incident History shall provide Save as Rule flow. & Selected incidents can prefill editable rule decisions. \\
\end{longtable}

\subsection{Reconnect, Submission, Auth, UI, Projector, and Installer}

\begin{longtable}{L{0.15\textwidth} L{0.47\textwidth} L{0.28\textwidth}}
\textbf{ID} & \textbf{Requirement} & \textbf{Acceptance Criteria} \\
\hline
FR-REC-001 & Disconnect shall be treated as network state, not local shutdown. & Monitors, GUI, recorder, and incident engine remain active. \\
FR-REC-002 & Disconnected incidents shall be persisted with stable sequence metadata. & Entries include sequence, queued time, and buffered flag. \\
FR-REC-003 & Buffered incidents and pending evidence shall flush after reconnect. & Unacknowledged items remain on disk until acknowledged. \\
FR-SUB-001 & The client shall build and upload final submission bundles. & Server validates UUID, size, file type, and checksum where supplied. \\
FR-AUTH-001 & The client shall ask the server for auth status before deciding preflight behavior. & Failure falls back to strict local auth. \\
FR-AUTH-002 & Disabled CATS or AD auth shall create admin validation state, not automatic admission. & Admin approves or denies through server commands. \\
FR-IPC-001 & Local IPC shall support loopback WebSocket with token authentication. & Non-loopback peers and invalid tokens are rejected. \\
FR-IPC-002 & Local IPC shall preserve stdio fallback. & Auto mode uses WebSocket only when IPC environment exists. \\
FR-UI-001 & Tk and Qt dashboards shall provide equivalent workflows. & Clients, incidents, process database, rules, settings, folders, and commands are available. \\
FR-UI-002 & Dashboard refresh shall preserve scroll, selection, focus, and horizontal scroll where possible. & Routine timer updates do not rebuild tables. \\
FR-PROJ-001 & The server shall serve a read-only projector page and SSE feed. & Public-safe aggregate state is streamed without controls. \\
FR-PROJ-002 & Projector assets shall be separated into HTML, CSS, and JavaScript. & Server serves separate static files. \\
FR-INS-001 & The offline installer shall install packages into a shared virtual environment rather than global site-packages. & Packages install under \texttt{C:/ProgramData/May\_04\_Deniz/python\_env}. \\
FR-INS-002 & The installer shall copy runnable app code to a controlled shared folder and verify bundle hashes when a manifest exists. & Launchers run from installed app path and \texttt{manifest.sha256} is checked. \\
\end{longtable}

\section{Non-Functional Requirements}

\begin{longtable}{L{0.18\textwidth} L{0.50\textwidth} L{0.22\textwidth}}
\textbf{ID} & \textbf{Requirement} & \textbf{Acceptance Criteria} \\
\hline
NFR-SEC-001 & Loopback IPC shall reject non-loopback peers. & Non-loopback addresses are rejected. \\
NFR-SEC-002 & Loopback IPC shall require a random token. & Missing or invalid token is rejected. \\
NFR-SEC-003 & Sensitive WebSocket events shall support signing and encryption. & Secured envelope contains timestamp, nonce, signature, and optional ciphertext. \\
NFR-SEC-004 & Projector payloads shall not expose sensitive data. & No login id, UUID, IP, process name, window title, path, or evidence field appears. \\
NFR-REL-001 & Reconnect shall not interrupt local monitoring. & Client logs keep growing during transient disconnect. \\
NFR-UX-001 & Dashboard live updates shall preserve interaction state. & Scrollbars and selections do not reset during routine updates. \\
NFR-UX-002 & Projector UI shall be readable on low-resolution projection. & Large type, high contrast, and minimal dense detail are used. \\
NFR-COMP-001 & Offline wheels shall match target Python ABI and platform. & Python 3.14 requires a rebuilt cp314 wheelhouse. \\
\end{longtable}

\section{External Interfaces}

\subsection{HTTP Routes}

\begin{longtable}{L{0.15\textwidth} L{0.30\textwidth} L{0.45\textwidth}}
\textbf{Method} & \textbf{Route} & \textbf{Purpose} \\
\hline
GET & \texttt{/health} & Server identity and availability check. \\
GET & \texttt{/auth/status?login\_id=<id>} & Auth requirement and validation status lookup. \\
POST & \texttt{/login} & Login and session UUID assignment. \\
GET & \texttt{/exam/config?id=<uuid>} & Exam configuration retrieval. \\
GET & \texttt{/exam/files?id=<uuid>} & Exam materials download. \\
POST & \texttt{/exam/submission?id=<uuid>} & Final submission upload. \\
POST & \texttt{/client/artifact?id=<uuid>} & Runtime artifact or evidence upload. \\
GET & \texttt{/ws?id=<uuid>} & LAN WebSocket runtime connection. \\
GET & \texttt{/projector} & Read-only projector page. \\
GET & \texttt{/projector/events} & Projection-safe SSE feed. \\
\end{longtable}

\subsection{Loopback IPC Channels}

\begin{longtable}{L{0.30\textwidth} L{0.25\textwidth} L{0.35\textwidth}}
\textbf{Channel} & \textbf{Direction} & \textbf{Purpose} \\
\hline
\texttt{manager.console\_command} & Manager to CLI & Sends console commands. \\
\texttt{server.dashboard\_state} & Server CLI to dashboard & Pushes state and command output. \\
\texttt{dashboard.command} & Dashboard to server CLI & Sends structured dashboard actions. \\
\texttt{client.timer\_state} & Client CLI to timer UI & Pushes timer, upload, and folder state. \\
\texttt{timer.command} & Timer UI to client CLI & Sends student UI actions. \\
\texttt{process.lifecycle} & Parent/child processes & Sends lifecycle notifications. \\
\end{longtable}

\section{Data Requirements}

\begin{longtable}{L{0.38\textwidth} L{0.18\textwidth} L{0.34\textwidth}}
\textbf{Path} & \textbf{Owner} & \textbf{Requirement} \\
\hline
\texttt{data/server/server\_users.json} & Server & Persist users and session state. \\
\texttt{data/server/incidents.jsonl} & Server & Append incident history. \\
\texttt{data/server/exam\_policy.json} & Server & Persist normalized exam policy. \\
\texttt{data/server/process\_definitions.json} & Server & Persist process definitions. \\
\texttt{data/server/incident\_rules.json} & Server & Persist incident rules. \\
\texttt{data/server/artifacts/*} & Server & Store evidence artifacts. \\
\texttt{data/server/submissions/*} & Server & Store final submissions. \\
\texttt{data/client/\{uuid\}/buffer/*} & Client & Persist buffered incidents and evidence retry state. \\
\texttt{data/client/\{uuid\}/recordings/*} & Client & Store replay cache and saved replays. \\
\texttt{Desktop/Exam/DD-MM-YYYY} & Client & Store extracted exam materials. \\
\end{longtable}

\section{Validation}

The expected automated validation commands are:

\begin{verbatim}
cd May_04_Deniz
python -m compileall -q .
python -m unittest discover -s tests
python -m pip check
\end{verbatim}

The latest local validation run while preparing this package passed compileall,
passed 168 tests, and reported no broken pip requirements.

\section{Traceability Summary}

\begin{longtable}{L{0.24\textwidth} L{0.28\textwidth} L{0.38\textwidth}}
\textbf{Area} & \textbf{Requirement IDs} & \textbf{Primary Modules} \\
\hline
Server runtime & FR-SRV-* & \texttt{server.main}, \texttt{server.app}, \texttt{server.handlers}, \texttt{server.tasks} \\
Client runtime & FR-CLI-* & \texttt{client.main}, \texttt{client.ws\_client}, \texttt{client.preflight} \\
Monitoring & FR-MON-* & \texttt{client.custommodules}, \texttt{client.incidents}, \texttt{common.process\_definitions} \\
Incident rules & FR-INC-* & \texttt{common.incident\_rules}, \texttt{server.state}, \texttt{server.ui} \\
Reconnect & FR-REC-* & \texttt{client.runtime}, \texttt{client.incident\_buffer}, \texttt{client.ws\_client} \\
Authentication & FR-AUTH-* & \texttt{server.auth\_validation}, \texttt{server.tasks}, \texttt{client.preflight} \\
IPC and UI & FR-IPC-*, FR-UI-* & \texttt{common.ipc\_ws}, \texttt{launcher\_ui}, \texttt{server.ui}, \texttt{client.ui} \\
Projector & FR-PROJ-* & \texttt{server.projector}, \texttt{server.static.projector} \\
Installer & FR-INS-* & \texttt{offline\_installer} \\
\end{longtable}

\end{document}
